% GENERATED FILE. Edit canonical lessons, then run npm run generate:books.
% canonical-source-hash: fnv1a64:52ee14457780d1f5
% canonical-lessons: HI-C174-lautana, HI-C174-sajana, HI-C174-sabhalna, HI-C174-sudharna, HI-C174-tolna

\chapter{To give back, To decorate, To look after, To improve, To weigh}
\label{ch:hi-to-give-back-to-decorate-to-look-after-to-improve-to-weigh}
\hlchaptermodality{174}

\begin{chapteropening}
\textbf{By the end of this chapter:} \emph{I can say to give back, to decorate, to look after, to improve and to weigh.}

Complete the last lesson of chapter 174: I can say to give back, to decorate, to look after, to improve and to weigh.
\end{chapteropening}

\section[\texorpdfstring{\hi{लौटाना} (lauṭānā)}{लौटाना (lauṭānā)}]{\hi{लौटाना} (lauṭānā) — to give back}

\label{lesson:HI-C174-lautana}



\begin{quote}
Before the new one: say the Hindi for to choose, then the Hindi for to ask for.
\end{quote}

\subsection*{You'll want to know: \hi{लौटाना}}
\textbf{\hi{लौटाना}} — \emph{lauṭānā} — \textquotedblleft{}to give back\textquotedblright{}.

\begin{cousinweb}
\textbf{\hi{लौटाना}} (\emph{lauṭānā}) is the causing form of \emph{lauṭnā}, \textquotedblleft{}to come back\textquotedblright{}: to give back is to make something return. The origin of \emph{lauṭnā} is uncertain, and CDIAL has no entry for it.
\end{cousinweb}

\begin{grammarlens}[title={What you've built}]
One more word for this chapter.
\end{grammarlens}

\subsection*{Your turn}
\begin{itemize}
\raggedright
  \item \emph{Say it:} \emph{lauṭānā}
  \item \emph{Say it:} \emph{lauṭānā}, once more
  \item \emph{Say it:} say \emph{lauṭānā}
  \item \emph{Recall:} say the Hindi for to choose, then the Hindi for to ask for, then say \emph{lauṭānā} again
\end{itemize}

\begin{tcolorbox}[breakable,colback=teal!4,colframe=teal!35!black,title={Before you move on}]
What does \hi{लौटाना} mean? (\textquotedblleft{}To give back\textquotedblright{}.) Say it once more.
\end{tcolorbox}
\section[\texorpdfstring{\hi{सजाना} (sajānā)}{सजाना (sajānā)}]{\hi{सजाना} (sajānā) — to decorate}

\label{lesson:HI-C174-sajana}



\begin{quote}
Before the new one: say the Hindi for to ask for, then the Hindi for to give back.
\end{quote}

\subsection*{You'll want to know: \hi{सजाना}}
\textbf{\hi{सजाना}} — \emph{sajānā} — \textquotedblleft{}to decorate\textquotedblright{}.

\begin{cousinweb}
\textbf{\hi{सजाना}} (\emph{sajānā}) is a tadbhava, inherited and reshaped by sound change. CDIAL lists it, as the causing form of \emph{sājnā}, \textquotedblleft{}to prepare, decorate\textquotedblright{}, under Sanskrit \emph{sajjayati}, \textquotedblleft{}equips, prepares\textquotedblright{} (CDIAL 13091).
\end{cousinweb}

\begin{grammarlens}[title={What you've built}]
One more word for this chapter.
\end{grammarlens}

\subsection*{Your turn}
\begin{itemize}
\raggedright
  \item \emph{Say it:} \emph{sajānā}
  \item \emph{Say it:} \emph{sajānā}, once more
  \item \emph{Say it:} say \emph{sajānā}
  \item \emph{Recall:} say the Hindi for to ask for, then the Hindi for to give back, then say \emph{sajānā} again
\end{itemize}

\begin{tcolorbox}[breakable,colback=teal!4,colframe=teal!35!black,title={Before you move on}]
What does \hi{सजाना} mean? (\textquotedblleft{}To decorate\textquotedblright{}.) Say it once more.
\end{tcolorbox}
\section[\texorpdfstring{\hi{सँभालना} (sãbhālnā)}{सँभालना (sãbhālnā)}]{\hi{सँभालना} (sãbhālnā) — to look after, to handle}

\label{lesson:HI-C174-sabhalna}



\begin{quote}
Before the new one: say the Hindi for to give back, then the Hindi for to decorate.
\end{quote}

\subsection*{You'll want to know: \hi{सँभालना}}
\textbf{\hi{सँभालना}} — \emph{sãbhālnā} — \textquotedblleft{}to look after, to handle\textquotedblright{}.

\begin{cousinweb}
\textbf{\hi{सँभालना}} (\emph{sãbhālnā}) is a tadbhava, a word inherited from Sanskrit and reshaped by sound change. CDIAL lists it in its entry for Sanskrit \emph{saṁbharati}, \textquotedblleft{}brings together, prepares\textquotedblright{} (CDIAL 12961), among a group of forms with \emph{-l-} that it places under neither of its two numbered headwords. The same entry puts Hindi \emph{sãbhārnā}, \textquotedblleft{}to sustain\textquotedblright{}, under the causing form \emph{saṁbhārayati}, \textquotedblleft{}maintains\textquotedblright{}.
\end{cousinweb}

\begin{grammarlens}[title={What you've built}]
One more word for this chapter.
\end{grammarlens}

\subsection*{Your turn}
\begin{itemize}
\raggedright
  \item \emph{Say it:} \emph{sãbhālnā}
  \item \emph{Say it:} \emph{sãbhālnā}, once more
  \item \emph{Say it:} say \emph{sãbhālnā}
  \item \emph{Recall:} say the Hindi for to give back, then the Hindi for to decorate, then say \emph{sãbhālnā} again
\end{itemize}

\begin{tcolorbox}[breakable,colback=teal!4,colframe=teal!35!black,title={Before you move on}]
What does \hi{सँभालना} mean? (\textquotedblleft{}To look after, to handle\textquotedblright{}.) Say it once more.
\end{tcolorbox}
\section[\texorpdfstring{\hi{सुधारना} (sudhārnā)}{सुधारना (sudhārnā)}]{\hi{सुधारना} (sudhārnā) — to improve, to correct}

\label{lesson:HI-C174-sudharna}



\begin{quote}
Before the new one: say the Hindi for to decorate, then the Hindi for to look after.
\end{quote}

\subsection*{You'll want to know: \hi{सुधारना}}
\textbf{\hi{सुधारना}} — \emph{sudhārnā} — \textquotedblleft{}to improve, to correct\textquotedblright{}.

\begin{cousinweb}
\textbf{\hi{सुधारना}} (\emph{sudhārnā}) is a tadbhava, a word inherited and reshaped by sound change. It comes from \emph{*śuddhakāra}, \textquotedblleft{}making clean\textquotedblright{}, a form CDIAL reconstructs from its descendants (CDIAL 12521). It is built from \emph{śuddha}, \textquotedblleft{}pure\textquotedblright{}, and \emph{kāra}, \textquotedblleft{}making\textquotedblright{}.
\end{cousinweb}

\begin{grammarlens}[title={What you've built}]
One more word for this chapter.
\end{grammarlens}

\subsection*{Your turn}
\begin{itemize}
\raggedright
  \item \emph{Say it:} \emph{sudhārnā}
  \item \emph{Say it:} \emph{sudhārnā}, once more
  \item \emph{Say it:} say \emph{sudhārnā}
  \item \emph{Recall:} say the Hindi for to decorate, then the Hindi for to look after, then say \emph{sudhārnā} again
\end{itemize}

\begin{tcolorbox}[breakable,colback=teal!4,colframe=teal!35!black,title={Before you move on}]
What does \hi{सुधारना} mean? (\textquotedblleft{}To improve, to correct\textquotedblright{}.) Say it once more.
\end{tcolorbox}
\section[\texorpdfstring{\hi{तोलना} (tolnā)}{तोलना (tolnā)}]{\hi{तोलना} (tolnā) — to weigh}

\label{lesson:HI-C174-tolna}



\begin{quote}
Before the new one: say the Hindi for to look after, then the Hindi for to improve.
\end{quote}

\subsection*{You'll want to know: \hi{तोलना}}
\textbf{\hi{तोलना}} — \emph{tolnā} — \textquotedblleft{}to weigh\textquotedblright{}.

\begin{cousinweb}
\textbf{\hi{तोलना}} (\emph{tolnā}) is a tadbhava, a word inherited from Sanskrit and reshaped by sound change. It comes from \emph{tolayati}, \textquotedblleft{}lifts, weighs, considers\textquotedblright{} (CDIAL 5979).
\end{cousinweb}

\begin{grammarlens}[title={What you've built}]
One more word for this chapter.
\end{grammarlens}

\subsection*{Your turn}
\begin{itemize}
\raggedright
  \item \emph{Say it:} \emph{tolnā}
  \item \emph{Say it:} \emph{tolnā}, once more
  \item \emph{Say it:} say \emph{tolnā}
  \item \emph{Recall:} say the Hindi for to look after, then the Hindi for to improve, then say \emph{tolnā} again
\end{itemize}

\begin{tcolorbox}[breakable,colback=teal!4,colframe=teal!35!black,title={Before you move on}]
What does \hi{तोलना} mean? (\textquotedblleft{}To weigh\textquotedblright{}.) Say it once more.
\end{tcolorbox}
